% Additive complete algebra proof; no preamble or bibliography.
% Parent bibliography keys: prog:DNE, prog:arrival02.
\section{The original residue-dual basis, every conductor chart,
and the exact action-defect rank}

\subsection{Original objects and scope}
This calculation proves equations (3)--(17) of the supplied
native-class continuation \cite{prog:arrival02} in the original DNE
coordinates \cite{prog:DNE}. It also evaluates the action-defect rank
for every selected set and constructs the exact transition from every
original row minor to the consecutive chart. All residue expansions and
matrix eliminations used below are proved here.

Keep the actual simple-quartet data and the actual full conductor:
\begin{equation}\tag{NCB1}\label{ncb:data}
\begin{gathered}
a=k+4,\quad q=(a+1)^2,\quad d=q'=(a-7)^2,\quad
v=\operatorname{ord}_0E_A,\quad h=q-v,\quad m=h-d>0,\\
\beta_{a,u,w}=\frac a2+(2u-a)\delta+i(2w-a)\gamma,
\quad 0\leq u,w\leq a,\quad 0<\delta<\tfrac12,\quad\gamma>2,\\
E_A(z)=\sum_{r,t=0}^{8}a_{rt}e^{\beta_{8,r,t}z},
\qquad \mu_j=E_A^{(j)}(0),\qquad
\mu_0=\cdots=\mu_{v-1}=0,\quad\mu_v\ne0 .
\end{gathered}
\end{equation}
Every actual period, branch, coefficient, and finite guard remains
unchanged. The lower roots are
$\beta'_{\alpha'}=\beta_{a-8,u',w'}$, with their original ordering.
There are $d$ lower roots and $q$ upper roots. All roots within either
grid are distinct: equality of real and imaginary parts forces equality
of both indices. Every upper root is nonzero since its real part is
at least $a(1/2-\delta)>0$. Set
\begin{equation}\tag{NCB2}\label{ncb:chi}
\chi(S)=\prod_{\alpha=1}^{q}(S-\beta_\alpha)
       =\sum_{\ell=0}^{q}c_\ell S^\ell,\quad c_q=1,\quad
E=\mathbb C[S]/(\chi),\quad
\mathbb V_{n,\alpha}=\beta_\alpha^n\quad(0\leq n<q).
\end{equation}
Represent elements of $E$ by their unique remainders of degree below
$q$. For an ordered subset $B$ of coordinate indices, $E_B$ inserts
coordinates in increasing index order; $E_B^{\mathsf T}$ extracts them.
These are ordinary transpose conventions, not Hermitian adjoints.

The literal DNE conductor coefficient map
$C_a^{\mathsf T}:\mathbb C^d\to\mathbb C^q$ is
\begin{equation}\tag{NCB3}\label{ncb:conductor}
\sum_\alpha(C_a^{\mathsf T}\eta)_\alpha e^{\beta_\alpha z}
=E_A(z)\sum_{\alpha'}\eta_{\alpha'}e^{\beta'_{\alpha'}z}.
\end{equation}
Its upper $(u,w)$, lower $(u',w')$ entry is
$a_{u-u',w-w'}$ when both differences lie in $\{0,\ldots,8\}$, and
zero otherwise, because
$\beta_{8,r,t}+\beta_{a-8,u',w'}=\beta_{a,u'+r,w'+t}$.
No divided-native scalar is inserted in this map.

The map $C_a^{\mathsf T}$ is injective. If its value is zero, the
product in \eqref{ncb:conductor} vanishes identically. Since
$\mu_v\ne0$, $E_A$ is nonzero on an open disc; the lower exponential
polynomial vanishes on that disc and hence identically. Its first $d$
derivatives at zero form the lower Vandermonde times $\eta$. The
determinant is
$\prod_{\alpha'<\beta'}(\beta'_{\beta'}-\beta'_{\alpha'})\ne0$,
so $\eta=0$.
To verify this determinant formula, its determinant is an alternating
polynomial in the $d$ root variables, hence divisible by every
$\beta'_{\beta'}-\beta'_{\alpha'}$. Its total degree is
$0+1+\cdots+(d-1)$, the degree of their product. The coefficient of
$(\beta'_1)^0(\beta'_2)^1\cdots(\beta'_d)^{d-1}$ is one in both
polynomials, so their constant quotient is one.

Let $P_v:\mathbb C^h\to\mathbb C^q$ insert $v$ leading zero
coordinates. The full product rule applied to \eqref{ncb:conductor}
gives
\begin{equation}\tag{NCB4}\label{ncb:moment-conductor}
\begin{gathered}
\mathbb VC_a^{\mathsf T}=P_vJ,\qquad
J=P_v^{\mathsf T}\mathbb VC_a^{\mathsf T},\\
J_{i,\alpha'}=\sum_{r=0}^{i}
\binom{v+i}{r}\mu_{v+i-r}(\beta'_{\alpha'})^r
\qquad(0\leq i<h).
\end{gathered}
\end{equation}
All derivatives below order $v$ vanish. At order $v+i$, the product
rule originally ranges over $0\leq r\leq v+i$; terms with $r>i$
vanish because their conductor derivative has order below $v$.
Thus the displayed truncation retains every nonzero conductor
derivative. Since $\mathbb V$ and $P_v$ are injective, $J$ has rank $d$.

\subsection{Complete residue-dual moment basis}
For each upper root let
$p_\alpha(S)=\chi(S)/((S-\beta_\alpha)\chi'(\beta_\alpha))$.
Define the residue isomorphism and its moment-coordinate version:
\begin{equation}\tag{NCB5}\label{ncb:residue}
\mathcal L:\mathbb C^q\longrightarrow E,\qquad
\mathcal L(\nu)=\sum_\alpha\chi'(\beta_\alpha)\nu_\alpha p_\alpha,
\qquad F=\mathcal L\mathbb V^{-1}.
\end{equation}
Evaluation at the roots gives
$(\mathcal L^{-1}f)_\alpha=f(\beta_\alpha)/\chi'(\beta_\alpha)$,
so $\mathcal L$ is invertible. The roots are simple and no derivative
factor vanishes. The formula for a simple-pole residue gives
\[
\sum_\alpha\operatorname{Res}_{S=\beta_\alpha}
\frac{\mathcal L(\nu)(S)P(S)}{\chi(S)}\,dS
=\sum_\alpha\nu_\alpha P(\beta_\alpha).
\]

Define, for $0\leq n<q$,
\begin{equation}\tag{NCB6}\label{ncb:pn}
\mathfrak p_n(S)=\sum_{\ell=n+1}^{q}c_\ell S^{\ell-n-1}.
\end{equation}
Their distinct monic degrees are $q-1-n$, so they form a basis of $E$.
Their complete residue duality is
\begin{equation}\tag{NCB7}\label{ncb:duality}
\sum_\alpha\operatorname{Res}_{S=\beta_\alpha}
\frac{\mathfrak p_n(S)S^j}{\chi(S)}\,dS=\delta_{nj}
\quad(0\leq j,n<q),\qquad Fe_n=\mathfrak p_n .
\end{equation}
To prove the first equality, the sum of finite residues equals the
$S^{-1}$ coefficient at large $S$: partial fractions consist of a
polynomial and terms $b_\alpha/(S-\beta_\alpha)$, and each such term
has both residue and $S^{-1}$ coefficient $b_\alpha$.
For $j<n$, the numerator has degree at most $q-2$, so this coefficient
is zero. For $j=n$, it has degree $q-1$ and leading coefficient one.
For $j>n$, the exact identity
$S^{n+1}\mathfrak p_n=\chi-\sum_{\ell=0}^{n}c_\ell S^\ell$
implies
\[
\frac{\mathfrak p_nS^j}{\chi}
=S^{j-n-1}
-\frac{S^{j-n-1}\sum_{\ell=0}^{n}c_\ell S^\ell}{\chi}.
\]
The first term is a polynomial, and the second numerator has degree
at most $j-1\leq q-2$. Both have zero $S^{-1}$ coefficient.
This proves the first equality. The pairings of $Fe_n$ with $S^j$
are $(\mathbb V\mathbb V^{-1}e_n)_j=\delta_{nj}$, so the proved
basis duality identifies $Fe_n$ with $\mathfrak p_n$.

Fix any actual $d$-row minor
$I\subset\{0,\ldots,h-1\}$ with $\det J_I\ne0$,
where $J_I=E_I^{\mathsf T}J$. Write $J(I)=I^c$ and
\begin{equation}\tag{NCB8}\label{ncb:native}
T_I=\{v+j:j\in J(I)\},\qquad
Z_I=\mathbb V^{-1}P_vE_{J(I)},\qquad
L_I=\mathcal LZ_I=FE_{T_I}.
\end{equation}
Hence $\operatorname{im}L_I$ is exactly the span of
$\{\mathfrak p_n:n\in T_I\}$ in the original selected coordinates.

For every upper root $\lambda$, the finite geometric sums give
\[
\sum_{n=0}^{q-1}\lambda^n\mathfrak p_n(S)
=\sum_{\ell=1}^{q}c_\ell
\frac{S^\ell-\lambda^\ell}{S-\lambda}
=\frac{\chi(S)-\chi(\lambda)}{S-\lambda}.
\]
Consequently the original eigenclass is
\begin{equation}\tag{NCB9}\label{ncb:eigenclass}
v_\lambda=\frac{\chi(S)}{(S-\lambda)\chi'(\lambda)}
=\sum_{n=0}^{q-1}t_n\mathfrak p_n,\qquad
t_n=\frac{\lambda^n}{\chi'(\lambda)},\qquad
v_\lambda\notin\operatorname{im}L_I.
\end{equation}
Every $t_n$ is nonzero because $\lambda\ne0$ and
$\chi'(\lambda)\ne0$. A proper coordinate subset of the displayed
basis cannot contain a vector whose every coordinate is nonzero.
This proves the last assertion for every original admitted minor.

\subsection{Complete low, conductor, and native decomposition}
For an arbitrary $f\in E$, put $t=F^{-1}f$ and
$t_{\geq v}=P_v^{\mathsf T}t$. Define maps with their full types:
\begin{equation}\tag{NCB10}\label{ncb:HQ}
\begin{gathered}
H_I=J_I^{-1}E_I^{\mathsf T}:\mathbb C^h\to\mathbb C^d,\qquad
Q_I=E_{J(I)}^{\mathsf T}-J_{J(I)}H_I:
       \mathbb C^h\to\mathbb C^m,\\
\eta_I=H_It_{\geq v},\qquad z_I=Q_It_{\geq v}.
\end{gathered}
\end{equation}
Checking first the $I$ rows and then their complement proves
\begin{equation}\tag{NCB11}\label{ncb:split}
\begin{gathered}
JH_I+E_{J(I)}Q_I=I_h,\qquad H_IJ=I_d,\qquad Q_IJ=0,\\
H_IE_{J(I)}=0,\qquad Q_IE_{J(I)}=I_m.
\end{gathered}
\end{equation}
Together with $FP_vJ=\mathcal LC_a^{\mathsf T}$ from
\eqref{ncb:moment-conductor}, these give the full identity
\begin{equation}\tag{NCB12}\label{ncb:decomposition}
f=\underbrace{\sum_{n<v}t_n\mathfrak p_n}_{f^{\rm low}}
 +\underbrace{\mathcal L(C_a^{\mathsf T}\eta_I)}_{f_I^{\rm cond}}
 +\underbrace{L_Iz_I}_{f_I^{\rm nat}} .
\end{equation}
These three spaces form a direct sum. The low space has arbitrary
first $v$ moment coordinates and zero remaining coordinates; both
other spaces have their first $v$ coordinates zero. On the remaining
coordinates, \eqref{ncb:split} proves
$\mathbb C^h=\operatorname{im}J\oplus\operatorname{im}E_{J(I)}$.
Thus the full basis map
\begin{equation}\tag{NCB13}\label{ncb:full-chart}
B_I=[\,FE_{\{0,\ldots,v-1\}},\ \mathcal LC_a^{\mathsf T},\ L_I\,]:
\mathbb C^v\oplus\mathbb C^d\oplus\mathbb C^m\longrightarrow E
\end{equation}
is invertible, also when its first block is empty ($v=0$).

For $f=v_\lambda$, all selected coordinates of $t_{\geq v}$ are
nonzero and $d\geq1$, so $\eta_I\ne0$. Injectivity of the full
conductor map and of $\mathcal L$ gives
$v_{\lambda,I}^{\rm cond}\ne0$. If $v>0$ then
$v_\lambda^{\rm low}\ne0$ as well. The direct sum just proved implies
\begin{equation}\tag{NCB14}\label{ncb:nonzero}
v_\lambda^{\rm low}+v_{\lambda,I}^{\rm cond}\ne0 .
\end{equation}
This is an algebraic direct sum, not an orthogonal splitting.

The exact relative quotient represented by the native section is
\begin{equation}\tag{NCB15}\label{ncb:relative}
\frac{E}
 {\operatorname{im}(FE_{\{0,\ldots,v-1\}})
          +\operatorname{im}(\mathcal LC_a^{\mathsf T})}
\ \xrightarrow[\ [f]\mapsto Q_IP_v^{\mathsf T}F^{-1}f\ ]{\ \cong\ }
\mathbb C^m .
\end{equation}
Its inverse is $z\mapsto[L_Iz]$. Its kernel and both inverse identities
follow from \eqref{ncb:split}. This proves the exact map relating the
full arithmetic class and its native representative; it retains their
proved nonzero difference in $E$.

\subsection{Every original minor to every other minor}
For another admitted $d$-row minor $K$, let
\begin{equation}\tag{NCB16}\label{ncb:transition}
D_{K\leftarrow I}=H_KE_{J(I)}:\mathbb C^m\to\mathbb C^d,\qquad
M_{K\leftarrow I}=Q_KE_{J(I)}:\mathbb C^m\to\mathbb C^m .
\end{equation}
Apply \eqref{ncb:split} for $K$ to $E_{J(I)}$ to obtain
\begin{equation}\tag{NCB17}\label{ncb:section-transition}
\begin{gathered}
E_{J(I)}=JD_{K\leftarrow I}
              +E_{J(K)}M_{K\leftarrow I},\\
L_I=\mathcal LC_a^{\mathsf T}D_{K\leftarrow I}
              +L_KM_{K\leftarrow I}.
\end{gathered}
\end{equation}
This is the full changed-representative identity with its conductor
correction. Apply $Q_I$ to its first line to obtain
$M_{I\leftarrow K}M_{K\leftarrow I}=I_m$, and exchange the minors
to obtain the reverse product. Therefore
\begin{equation}\tag{NCB18}\label{ncb:inverse}
M_{K\leftarrow I}^{-1}=M_{I\leftarrow K},\qquad
\eta_K=\eta_I+D_{K\leftarrow I}z_I,\quad
z_K=M_{K\leftarrow I}z_I,\quad
f_K^{\rm low}=f_I^{\rm low}.
\end{equation}
The component identities follow by applying $H_K,Q_K$ to
$t_{\geq v}=J\eta_I+E_{J(I)}z_I$. In particular
\[
f_I^{\rm nat}=f_K^{\rm nat}
       +\mathcal LC_a^{\mathsf T}D_{K\leftarrow I}z_I,\qquad
f_K^{\rm cond}=f_I^{\rm cond}
       +\mathcal LC_a^{\mathsf T}D_{K\leftarrow I}z_I.
\]
Their total is unchanged because the same correction occurs on
opposite sides when the decompositions are compared.

For a third admitted minor $H$, applying $Q_H,H_H$ to
\eqref{ncb:section-transition} proves the oriented composition laws
\begin{equation}\tag{NCB19}\label{ncb:cocycle}
\begin{gathered}
M_{H\leftarrow I}=M_{H\leftarrow K}M_{K\leftarrow I},\\
D_{H\leftarrow I}
 =D_{K\leftarrow I}+D_{H\leftarrow K}M_{K\leftarrow I}.
\end{gathered}
\end{equation}

The exact determinant including ordering signs is also explicit.
For $I=(i_0<\cdots<i_{d-1})$ put
$\epsilon_I=(-1)^{\sum_{r=0}^{d-1}(i_r-r)}$.
The row permutation taking full increasing order to $(I,J(I))$
has this sign, since selected row $i_r$ crosses $i_r-r$ earlier
unselected rows. Under that permutation $W_I=[J,E_{J(I)}]$ is
block lower triangular with diagonal blocks $J_I,I_m$, so
$\det W_I=\epsilon_I\det J_I$. The first line of
\eqref{ncb:section-transition} gives
$W_I=W_K\left[\begin{smallmatrix}I_d&D\\0&M\end{smallmatrix}\right]$.
Consequently
\begin{equation}\tag{NCB20}\label{ncb:dettransition}
\det M_{K\leftarrow I}
 =\frac{\epsilon_I\det J_I}{\epsilon_K\det J_K}.
\end{equation}

The full coordinate transformation, including the unchanged low
coordinates, is
\begin{equation}\tag{NCB21}\label{ncb:full-transition}
T_{K\leftarrow I}=
\begin{pmatrix}I_v&0&0\\0&I_d&D_{K\leftarrow I}\\
0&0&M_{K\leftarrow I}\end{pmatrix},
\qquad B_I=B_KT_{K\leftarrow I}.
\end{equation}
For every original arithmetic Hermitian metric $G_N$ and the full
operator $A$ of multiplication by $S$ modulo $\chi$, set
$G_{I,N}=B_I^*G_NB_I$ and $A_I=B_I^{-1}AB_I$. Substitution yields
\begin{equation}\tag{NCB22}\label{ncb:gram-action}
G_{I,N}=T_{K\leftarrow I}^*G_{K,N}T_{K\leftarrow I},\qquad
A_I=T_{K\leftarrow I}^{-1}A_KT_{K\leftarrow I}.
\end{equation}
All nine Gram blocks, including all six off-diagonal pairings, are
retained. Multiplication of these identities gives
\begin{equation}\tag{NCB23}\label{ncb:centered-action}
\begin{split}
A_I^*G_{I,N}+G_{I,N}A_I-aG_{I,N}
={}&T_{K\leftarrow I}^*
(A_K^*G_{K,N}+G_{K,N}A_K-aG_{K,N})T_{K\leftarrow I}.
\end{split}
\end{equation}
Thus the full arithmetic action is preserved under the original
coordinate transition with its explicit conductor corrections.

\subsection{Consecutive chart with all conductor derivatives}
Let $I_0=\{0,\ldots,d-1\}$. Equation \eqref{ncb:moment-conductor}
gives the exact factorization
\begin{equation}\tag{NCB24}\label{ncb:triangular}
J_{I_0}=\mathsf T\mathbb V',\qquad
\mathsf T_{i,r}=
\begin{cases}\binom{v+i}{r}\mu_{v+i-r},&0\leq r\leq i<d,\\
0,&r>i,\end{cases}
\quad \mathbb V'_{r,\alpha'}=(\beta'_{\alpha'})^r .
\end{equation}
The lower-triangular diagonal entries are
$\binom{v+i}{i}\mu_v$. Taking the determinant of the actual product
therefore proves
\begin{equation}\tag{NCB25}\label{ncb:consecutive-det}
\det J_{I_0}
=\mu_v^d\left(\prod_{i=0}^{d-1}\binom{v+i}{i}\right)
 \prod_{\alpha'<\beta'}(\beta'_{\beta'}-\beta'_{\alpha'})\ne0 .
\end{equation}
All $\mu_j$ with $j>v$ remain in the matrix. Their absence from its
determinant is the proved effect of triangularity.

Its inverse retains every such derivative. For $x\in\mathbb C^d$,
solve $\mathsf T b=x$ by the finite recursion
\begin{equation}\tag{NCB26}\label{ncb:triangular-inverse}
b_0=\frac{x_0}{\mu_v},\qquad
b_i=\frac{x_i-\sum_{r=0}^{i-1}
          \binom{v+i}{r}\mu_{v+i-r}b_r}
          {\binom{v+i}{i}\mu_v}\quad(1\leq i<d).
\end{equation}
Write
$\chi_{\rm low}(S)=\prod_{\alpha'}(S-\beta'_{\alpha'})
=\sum_{\ell=0}^{d}c_\ell^{\rm low}S^\ell$
and
$\mathfrak p_r^{\rm low}(S)
=\sum_{\ell=r+1}^{d}c_\ell^{\rm low}S^{\ell-r-1}$.
The already proved residue-dual basis applied to $\chi_{\rm low}$
evaluates the inverse lower Vandermonde, giving
\begin{equation}\tag{NCB27}\label{ncb:full-inverse}
(J_{I_0}^{-1}x)_{\alpha'}
=\frac{\sum_{r=0}^{d-1}b_r\mathfrak p_r^{\rm low}(\beta'_{\alpha'})}
          {\chi_{\rm low}'(\beta'_{\alpha'})}.
\end{equation}
Indeed the vector on the right has lower moments $b$ by
\eqref{ncb:duality}, and those moments satisfy $\mathsf T b=x$.

For an arbitrary original chart $I$ and its native vector $z$, take
$x=E_{I_0}^{\mathsf T}E_{J(I)}z$ in
\eqref{ncb:triangular-inverse}--\eqref{ncb:full-inverse}. Their result
is exactly $D_{I_0\leftarrow I}z$. The native coordinates in the
consecutive chart are, for $d\leq j<h$,
\begin{equation}\tag{NCB28}\label{ncb:explicit-transition}
\begin{split}
(M_{I_0\leftarrow I}z)_{j-d}
={}&(E_{J(I)}z)_j\\
&-\sum_{\alpha'}\sum_{r=0}^{j}
\binom{v+j}{r}\mu_{v+j-r}(\beta'_{\alpha'})^r
                    (D_{I_0\leftarrow I}z)_{\alpha'} .
\end{split}
\end{equation}
This evaluates the chart transition in original coefficients, roots,
and derivatives. For the full eigenclass one instead sets
$x_i=\lambda^{v+i}/\chi'(\lambda)$ to calculate $\eta_{I_0}$,
and then sets
$z_{I_0,j-d}=t_{v+j}-(J\eta_{I_0})_j$.

The selected consecutive set is
$T_{I_0}=\{v+d,\ldots,q-1\}=\{q-m,\ldots,q-1\}$.
Its residue-dual basis elements have all monic degrees
$m-1,\ldots,0$, so their triangular coefficient matrix proves
\begin{equation}\tag{NCB29}\label{ncb:consecutive-image}
\operatorname{im}L_{I_0}=\mathbb C[S]_{\leq m-1}\subset E .
\end{equation}

\subsection{Action defect and exact rank for every selected set}
Multiplication by $S$ modulo the complete $\chi$ satisfies
\begin{equation}\tag{NCB30}\label{ncb:action}
A\mathfrak p_0=-c_0\mathfrak p_{q-1},\qquad
A\mathfrak p_n=\mathfrak p_{n-1}-c_n\mathfrak p_{q-1}
\quad(1\leq n<q).
\end{equation}
For the first equality use $S\mathfrak p_0=\chi-c_0$;
for the second use the literal polynomial identity
$S\mathfrak p_n=\mathfrak p_{n-1}-c_n$.
Here $\mathfrak p_{q-1}=1$.
For any selected $T\subset\{0,\ldots,q-1\}$, the omitted-coordinate
action is therefore
\begin{equation}\tag{NCB31}\label{ncb:RT}
\begin{gathered}
R_T=E_{T^c}^{\mathsf T}F^{-1}AFE_T,\\
(R_T)_{\ell n}
=\mathbf1_{\{n\geq1,\ \ell=n-1\}}
       -c_n\mathbf1_{\{\ell=q-1\}}
\qquad(\ell\in T^c,\ n\in T).
\end{gathered}
\end{equation}

Define the descending starts and their number by
$B(T)=\{n\in T:n\geq1,\ n-1\notin T\}$ and $b(T)=|B(T)|$.
The exact rank, including nongeneric vanishing coefficients, is
\begin{equation}\tag{NCB32}\label{ncb:rank}
\begin{gathered}
\operatorname{rank}R_T=b(T)+\varepsilon(T),\\
\varepsilon(T)=
\begin{cases}
1,&q-1\notin T\ \text{and some }n\in T\setminus B(T)
                                    \text{ has }c_n\ne0,\\
0,&\text{otherwise}.
\end{cases}
\end{gathered}
\end{equation}
For each $n\in B(T)$, row $n-1$ has exactly a single one in column
$n$. These $b(T)$ rows are independent and none is row $q-1$.
All other rows vanish except possibly row $q-1$, present exactly when
$q-1\notin T$. That row is $(-c_n)_{n\in T}$.
Subtracting its coordinates on columns $B(T)$ using the independent
rows leaves its coordinates on $T\setminus B(T)$; these add one rank
precisely in the case displayed. This proves \eqref{ncb:rank}.

For every nonempty proper $T$, the rank is positive. If $B(T)$ is
nonempty, this follows immediately. Otherwise descending from any
member of $T$ cannot leave $T$ before reaching zero, so $T$ is an
initial interval containing zero. Properness implies $q-1\notin T$,
and $c_0=(-1)^q\prod_\alpha\beta_\alpha\ne0$ gives
$\varepsilon(T)=1$. For the consecutive set $T_{I_0}$ there is just
one descending start $q-m=v+d>0$, and $q-1\in T_{I_0}$.
Thus its rank is exactly one.

\subsection{Original metric residual and its full Gram}
Let $G_N$ be any original positive Hermitian arithmetic quotient metric
on $E$, fix $T=T_I$, and set $L=FE_T$ and $\Gamma=L^*G_NL$.
The metric compression and residual, both with domain $\mathbb C^m$,
are
\begin{equation}\tag{NCB33}\label{ncb:compression}
A_{\rm nat,N}=\Gamma^{-1}L^*G_NAL,\qquad
\mathcal R_N=AL-LA_{\rm nat,N}.
\end{equation}
Write $\widetilde G=F^*G_NF$, and use $C=T^c$ solely as a block
index in the next formulas. Set
\begin{equation}\tag{NCB34}\label{ncb:lift}
\begin{gathered}
W_G=F(E_C-E_T\widetilde G_{TT}^{-1}\widetilde G_{TC}),\\
S_G=\widetilde G_{CC}
 -\widetilde G_{CT}\widetilde G_{TT}^{-1}\widetilde G_{TC}.
\end{gathered}
\end{equation}
This leaves the original $G_N$ unchanged. Direct multiplication gives
$L^*G_NW_G=0$ and $E_C^{\mathsf T}F^{-1}W_G=I_{q-m}$.
Hence $W_G$ is an injective map onto the $G_N$-orthogonal complement
of $\operatorname{im}L$, whose dimension is $q-m$; its inverse there
is literal omitted-coordinate extraction.
By \eqref{ncb:compression}, $\mathcal R_N$ has image in that
complement and omitted coordinates $R_T$. Uniqueness of this exact
lift proves
\begin{equation}\tag{NCB35}\label{ncb:residual}
\mathcal R_N=W_GR_T,\qquad
\mathcal R_N^*G_N\mathcal R_N=R_T^*S_GR_T,\qquad
\operatorname{rank}\mathcal R_N=\operatorname{rank}R_T .
\end{equation}
For the Gram identity, multiply \eqref{ncb:lift} to obtain
$W_G^*G_NW_G=S_G$, retaining its cross terms.
It is positive definite since $W_G$ is injective and $G_N$ is
positive definite. This proves the residual rank at every original
cutoff, with its exact value \eqref{ncb:rank}, and proves nonvanishing
for every original nonempty proper chart. No size estimate follows
from rank alone.

\subsection{What this finite calculation establishes}
The residue-dual basis, numerator image, every eigenclass
decomposition, nonzero omitted component, original companion action,
and complete consecutive determinant prove equations (3)--(17) of
\cite{prog:arrival02}. Equations
\eqref{ncb:transition}--\eqref{ncb:centered-action} provide the exact
comparison between every original minor and the consecutive chart.
The derivative recursion \eqref{ncb:triangular-inverse} and
\eqref{ncb:full-inverse} evaluate this comparison without deleting
conductor derivatives. The rank formula \eqref{ncb:rank} and residual
Gram \eqref{ncb:residual} strengthen the finite action calculation.

Under the programme's injective arithmetic realization, the nonzero
omitted vector \eqref{ncb:nonzero} remains nonzero. Its map from the
full remainder space to the native relative quotient is precisely
\eqref{ncb:relative}; changing the section is precisely
\eqref{ncb:section-transition}. The arithmetic pairing of the full
class consequently uses the complete block congruence
\eqref{ncb:gram-action} and \eqref{ncb:centered-action}.
This calculation does not identify the native section with the
original observation kernel or evaluate an absolute allocation or
individual projected sign.
